\documentclass[aspectratio=169]{beamer}
\usepackage{slaif}

\SlaifSetPresentationDate{30. maj 2026}
\SlaifSetPresentationFooter{Hitri seminar: sodobna UI v pedagoškem delu}

\begin{document}

\SlaifTitleSlide{
  title       = {Hitri seminar:\\Sodobna UI kot pomoč pri napornih\\pedagoških opravilih},
  author      = {Janez Perš},
  institution = {UL FE},
  event       = {Za pedagoške kolege in asistente},
  subtitle    = {Praktični delovni tokovi}
}

\SlaifMotivationSlide{
  title = {Zakaj je to posebej pomembno za mlajše sodelavce?},
  body  = {{\bfseries Mlajši sodelavci pogosto nosijo največji del operativnega pedagoškega bremena.}\\[8pt]
  Oddaje, vaje, napake, povratne informacije, zagovori in popravljanje zahtevajo strokovno presojo, a se pogosto ponavljajo.\\[8pt]
  Današnja teza: UI je koristna tam, kjer ohrani odgovornost pri učitelju, zmanjša ponavljanje in poveča preglednost dela.}
}

\NoTextSlide{
  title = {Ocenjevanje ni drobna administracija},
  subtitle = {Je strokovno delo, ki se ponavlja, se kopiči in tiho poje čas za boljše poučevanje.},
  body = {
    \SlaifStatisticPanels{
      \SlaifStatistic{9,9}{ur / teden}{Povprečen čas, ki ga učitelji v anketi porabijo za ocenjevanje in označevanje.}
      \SlaifStatistic{--}{energija}{Obsežno ocenjevanje porablja pozornost, ki bi lahko šla v načrtovanje predmeta.}
      \SlaifStatistic{!}{čustveni strošek}{Na univerzi je ocenjevanje doživeto manj pozitivno kot raziskovanje in poučevanje.}
    }
    \SlaifFootnotePanel[fill=SlaifPanelBlue]{Learnosity (2025): učitelji poročajo o 9,9 h ocenjevanja na teden (n=258).\\Schwab et al. (2024): ocenjevanje je čustveno manj pozitivno kot raziskovanje in poučevanje.\\Terada in Merrill (2024): pogosto ocenjevanje jemlje čas in energijo za izboljševanje pouka.}
  }
}

\NoTextSlide{
  title = {V STEM predmetih je breme še težje},
  subtitle = {Ne preverjamo samo končnega odgovora, ampak pot, skice, formule, napake in razloge.},
  body = {
    \ColoredPanels{
      \ColoredPanel{Več korakov}{}{\SlaifPanelCenteredItems{\SlaifPanelTextItem{Študent lahko pride do napačnega rezultata z dobro metodo ali do prave številke po slabi poti.}}}{}
      \ColoredPanel{Več modalnosti}{}{\SlaifPanelCenteredItems{\SlaifPanelTextItem{Rokopis, diagrami, vezja, enačbe in prosti komentarji so del odgovora, ne okras.}}}{}
      \ColoredPanel{Več zakasnitve}{}{\SlaifPanelCenteredItems{\SlaifPanelTextItem{Ko se povratna informacija zavleče, je manj uporabna za učenje in naslednje vaje.}}}{}
    }
    \SlaifFootnotePanel[fill=SlaifPanelBlue]{Tan et al. (2025): STEM ocenjevanje vključuje večkorakovne rešitve, formule, diagrame in grafične artefakte.\\Tan et al. (2025): pri večjih obremenitvah je pravočasna povratna informacija poseben problem.\\Singh et al. (2017): digitalizacija izboljša hitrost in konsistentnost, presoja pa ostane človeško delo.}
  }
}

\NoTextSlide{
  title = {Teza seminarja: razbremeniti ponavljanje, ne presoje},
  subtitle = {UI je uporabna šele, ko je delo ponovljivo, preverljivo in popravljivo.},
  body = {
    \ColoredPanels{
      \ColoredPanel{Priprava}{iz gradiv}{\SlaifPanelCenteredItems{\SlaifPanelTextItem{{\bfseries Hitrejši prvi osnutek}}\SlaifPanelTextItem{Viri, slike, struktura in poudarki se uredijo v pregledno izhodišče.}}}{}
      \ColoredPanel{Preverjanje}{oddaj}{\SlaifPanelCenteredItems{\SlaifPanelTextItem{{\bfseries Boljša vprašanja}}\SlaifPanelTextItem{Sistem pripravi izhodišča za pogovor, učitelj pa vodi presojo.}}}{}
      \ColoredPanel{Ocenjevanje}{STEM odgovorov}{\SlaifPanelCenteredItems{\SlaifPanelTextItem{{\bfseries Več nadzora}}\SlaifPanelTextItem{Reference, pravila, več ocenjevalcev in ročni pregled zmanjšajo slepe odločitve.}}}{}
      \ColoredPanel{Dokumentiranje}{odločitev}{\SlaifPanelCenteredItems{\SlaifPanelTextItem{{\bfseries Več sledljivosti}}\SlaifPanelTextItem{Vsaka trditev mora imeti vir, razlog in možnost popravka.}}}{}
    }
    \SlaifFootnotePanel[fill=SlaifPanelMint]{Singh et al. (2017): urejen digitalni delovni tok zmanjša režijo in izboljša konsistentnost.\\Tan et al. (2025): avtomatizacija je smiselna tam, kjer zmanjša obremenitev in pospeši povratno informacijo.\\Schwab et al. (2024): zmanjševati je treba tudi čustveni strošek ocenjevanja, ne samo število klikov.}
  }
}

\NoTextSlide{
  title = {Meja je jasna: orodje ne sme prevzeti odgovornosti},
  subtitle = {Najslabša uporaba UI je tista, ki pospeši delo tako, da skrije presojo.},
  body = {
    \ColoredPanels{
      \ColoredPanel{Ne}{}{\SlaifPanelCenteredItems{\SlaifPanelTextItem{UI ne pozna lokalnih pravil, ciljev predmeta in mejnih primerov sama od sebe.}}}{}
      \ColoredPanel{Ne}{}{\SlaifPanelCenteredItems{\SlaifPanelTextItem{Modelov odgovor ni odločba; je gradivo za presojo, preverjanje in izboljšavo.}}}{}
      \ColoredPanel{Ne}{}{\SlaifPanelCenteredItems{\SlaifPanelTextItem{Detektorji niso stabilna pedagoška podlaga za obtoževanje študentov.}}}{}
      \ColoredPanel{Ne}{}{\SlaifPanelCenteredItems{\SlaifPanelTextItem{Če postopka ni mogoče razložiti študentu, ni zrel za pedagoško rabo.}}}{}
    }
    \SlaifFootnotePanel[fill=SlaifPanelBlue]{Singh et al. (2017): uspešni sistemi najprej uredijo človeški delovni tok.\\Tan et al. (2025): avtomatizacija naj zmanjša obremenitev in pospeši povratno informacijo, ne odstrani odgovornosti.\\Schwab et al. (2024): zmanjšanje bremena ne sme razvrednotiti strokovne presoje.}
  }
}

\SlaifGreenDividerSlide
  {Primer 1}
  {UI-podprta priprava prosojnic za STEM}
  {Od gradiv do preverljivega osnutka, ki ga učitelj lahko popravi.}

\NoTextSlide{
  title = {Postopek priprave prosojnic iz poljubnih gradiv},
  subtitle = {Repozitorij vsebuje navodila agentu ter spretnosti za uvoz virov, izdelavo predstavitev in razvoj predlog LaTeX.},
  body = {
    \ColoredPanels{
      \ColoredPanel{1. Viri}{}{\SlaifPanelCenteredItems{\SlaifPanelTextItem{skripta\\članek\\stare prosojnice\\slike in tabele}}}{}
      \ColoredPanel{2. Uvoz}{}{\SlaifPanelCenteredItems{\SlaifPanelTextItem{inventar\\izrezi\\struktura\\viri}}}{}
      \ColoredPanel{3. Zasnova}{}{\SlaifPanelCenteredItems{\SlaifPanelTextItem{zaporedje\\didaktični poudarki\\vizualni ritem}}}{}
      \ColoredPanel{4. Pregled}{}{\SlaifPanelCenteredItems{\SlaifPanelTextItem{prevod PDF\\PNG pregled\\popravki}}}{}
    }
    \SlaifPlainBulletBlock{
      \item GitHub repozitorij: \href{https://github.com/ulfe-lmi/slaif-presentation-builder}{\underline{https://github.com/ulfe-lmi/slaif-presentation-builder}}
      \item Agent iterira izgled, dokler izgled ni v redu.
      \item Učitelj ostane avtor: določi raven razlage, izbriše odvečno, popravi poudarke in potrdi končno sporočilo.
      \item Učitelj komunicira z agentom; ni treba, da se sploh zaveda, da je v ozadju LaTeX.
    }
  }
}

\NoTextSlide{
  title = {Uvoz vira: kaj mora agent oddati},
  subtitle = {Dober uvoz ni povzetek, ampak mapa dokazov za kasnejšo predstavitev.},
  body = {
    \SlaifStatisticPanels[y=128,totalWidth=820,x=70,height=155]{
      \SlaifStatistic{1}{Struktura}{kaj je v virih in kako bi to lahko postalo zaporedje prosojnic}
      \SlaifStatistic{2}{Gradniki}{slike, izrezi, tabele, strani in drugi elementi, ki jih je smiselno uporabiti}
      \SlaifStatistic{3}{Sledljivost}{od kod prihaja trditev ali vizualni element, da učitelj lahko preveri vir}
    }
    \SlaifPlainBulletBlock[x=90,y=350,width=820]{
      \item To ni samo formalnost.
      \item Dobre prosojnice so kompromis med strukturo in izgledom.
      \item Rezultat "ChatGPT, naredi mi PPTX" je lep, ampak nekonsistenten izdelek - popolnoma brez strukture.
      \item Struktura koristi, ko hočemo gradivo ponovno uporabiti ali popraviti.
    }
  }
}

\NoTextSlide{
  title = {Ustvarjanje predstavitve: kaj mora agent preveriti},
  subtitle = {Uspešen prevod LaTeX-a še ne pomeni uporabne predstavitve.},
  body = {
    \ColoredPanels{
      \ColoredPanel{Vsebina}{}{\SlaifPanelCenteredItems{\SlaifPanelTextItem{slovenščina\\jasno zaporedje\\primerna gostota}}}{}
      \ColoredPanel{Vizualno}{}{\SlaifPanelCenteredItems{\SlaifPanelTextItem{berljive slike\\pravilni izrezi\\dovolj praznega prostora}}}{}
      \ColoredPanel{Prevod}{}{\SlaifPanelCenteredItems{\SlaifPanelTextItem{lokalni PDF\\brez napak\\ponovljiv ukaz}}}{}
      \ColoredPanel{Pregled}{}{\SlaifPanelCenteredItems{\SlaifPanelTextItem{PNG izris\\ročno gledanje\\popravki}}}{}
    }
    \SlaifCommentaryPanel{Končni standard}{Agent mora sam iskati prekrivanje besedila, premajhne črke, slabe izreze in zlomljene postavitve.}
  }
}

\SlaifBlueDividerSlide
  {Primer 2}
  {Priprava vprašanj za zagovor samostojno izvedenih aktivnosti}
  {Kako preverimo, ali študenti razumejo oddane izdelke.}

\NoTextSlide{
  title = {Študenti uporabljajo umetno inteligenco},
  subtitle = {Zanesljive detekcije uporabe UI ni, zato mora zagovor preverjati razumevanje konkretne oddaje.},
  body = {
    \ColoredPanels{
      \ColoredPanel{Slaba smer}{}{\SlaifPanelCenteredItems[y=236,rowStep=54,itemHeight=42]{\SlaifPanelTextItem{{\bfseries "Ali si uporabil UI?"}}\SlaifPanelTextItem{Nezanesljivo, težko dokazljivo in pogosto pedagoško neproduktivno.}}}{}
      \ColoredPanel{Dobra smer}{}{\SlaifPanelCenteredItems[y=236,rowStep=54,itemHeight=42]{\SlaifPanelTextItem{{\bfseries "Ali razumeš to, kar si oddal?"}}\SlaifPanelTextItem{Vprašanja o kodi, grafu, parametru, odločitvi, napaki ali nenavadni razlagi.}}}{}
    }
    \SlaifPlainBulletBlock[variant=blue,x=90,y=372,width=790]{
      \item GitHub repozitorij: \href{https://github.com/ulfe-lmi/slaif-defense-grading}{\underline{https://github.com/ulfe-lmi/slaif-defense-grading}}
      \item Ne dokazujemo uporabe UI, ampak preverjamo razumevanje oddanega izdelka.
      \item Slab zagovor je pedagoško pomemben signal, ne glede na vzrok.
      \item UI pripravi ciljana vprašanja, ne sklepa o avtorstvu.
    }
  }
}

\NoTextSlide{
  title = {Delovni tok za zagovore vaj},
  subtitle = {Oddaje se izvozijo iz Moodla, psevdonimizirajo, primerjajo N x N in pretvorijo v vprašanja za zagovor.},
  body = {
    \ColoredPanels{
      \ColoredPanel{1. Moodle}{}{\SlaifPanelCenteredItems{\SlaifPanelTextItem{izvoz oddaj\\seznam študentov\\navodilo}}}{}
      \ColoredPanel{2. Priprava}{}{\SlaifPanelCenteredItems{\SlaifPanelTextItem{psevdonimizacija\\besedila\\koda\\priloge}}}{}
      \ColoredPanel{3. Primerjava}{}{\SlaifPanelCenteredItems{\SlaifPanelTextItem{N x N oddaj\\podobnosti\\razlike}}}{}
      \ColoredPanel{4. Vprašanja}{}{\SlaifPanelCenteredItems{\SlaifPanelTextItem{vidiki\\razlogi\\pričakovani odgovori}}}{}
    }
    \SlaifPlainBulletBlock[x=90,y=370,width=790]{
      \item Agent sistematično izvede N x N primerjav oddanih izdelkov.
      \item Izhod niso sodbe, ampak podlaga za ciljana vprašanja.
      \item Vsako vprašanje vsebuje razlog, vprašanje študentu in pričakovani odgovor.
    }
  }
}

\NoTextSlide{
  title = {Kakšno vprašanje je uporabno na zagovoru?},
  subtitle = {Dobro vprašanje je specifično za študenta, preverljivo in vezano na njegovo oddajo.},
  body = {
    \ColoredPanels{
      \ColoredPanel{Preširoko}{}{\SlaifPanelCenteredItems[y=236,rowStep=43,itemHeight=31]{\SlaifPanelTextItem{Razloži regresijo.}\SlaifPanelTextItem{Ali razumeš program?}\SlaifPanelTextItem{Kaj je rezultat?}}}{}
      \ColoredPanel{Uporabno}{}{\SlaifPanelCenteredItems[y=236,rowStep=43,itemHeight=31]{\SlaifPanelTextItem{Zakaj si tukaj izbral ta prag?}\SlaifPanelTextItem{Kaj se zgodi, če zamenjamo parameter?}\SlaifPanelTextItem{Zakaj je ta graf drugačen od pričakovanega?}}}{}
    }
    \SlaifPlainBulletBlock[variant=blue,x=90,y=382,width=790]{
      \item Vprašanja so specifična za posameznega študenta in njegovo konkretno oddajo.
      \item Kjer je mogoče, se vprašanja ne ponavljajo.
      \item Standardni naučeni odgovori izgubijo vrednost.
      \item Pedagog dobi individualna vprašanja brez ročne priprave za vsakega študenta.
    }
  }
}

\SlaifGreenDividerSlide
  {Primer 3}
  {UI-podprto ocenjevanje STEM izpitov}
  {Do GPT-5.2 je bil to doslej nerešen problem.}

\NoTextSlide{
  title = {Ročno pisani izpiti v tehniki in naravoslovju},
  subtitle = {Ocenjevanje mora biti iz perspektive študenta 100\% kompatibilno z dosedanjim načinom pisnih izpitov.},
  body = {
    \ColoredPanels{
      \ColoredPanel{Vhod je kompleksen}{}{\SlaifPanelCenteredItems{\SlaifPanelTextItem{ročno pisanje\\enačbe\\diagrami\\popravki}}}{}
      \ColoredPanel{Ocena je postopna}{}{\SlaifPanelCenteredItems{\SlaifPanelTextItem{vmesni koraki\\točke po pravilih\\razlaga napak}}}{}
      \ColoredPanel{OCR ni dovolj}{}{\SlaifPanelCenteredItems{\SlaifPanelTextItem{matematika\\skice\\nečitljivost\\kontekst}}}{}
      \ColoredPanel{Zato: sistem}{}{\SlaifPanelCenteredItems{\SlaifPanelTextItem{referenca\\pravila\\več ocen\\nadzor}}}{}
    }
    \SlaifCommentaryPanel{Zakaj je to dober primer}{Pri ročno pisanih izpitih "samo uporabi model" hitro odpove; urejen cevovod pa lahko pomaga, ker loči prepoznavanje, pravila, ocenjevanje, validacijo in ročni pregled.}
  }
}

\TwoColumnNoTextSlide{
  title = {Zakaj je STEM rokopis drugačen problem},
  subtitle = {Sken pisnega izpita je zelo kompleksen ročno narejen dokument.},
  left = {
    \SlaifPlainBulletBlock[variant=blue,x=72,y=178,width=390,height=260]{
      \item besedilo in formule sta prepletena;
      \item diagram je del odgovora;
      \item nekaj je prečrtano ali popravljeno;
      \item ocena mora razložiti delne točke.
    }
  },
  right = {
    \placeimage{585}{145}{315}{319}{stem-handwritten-example-nocaption.png}
    \SlaifFootnotePanel[fill=SlaifPanelBlue]{Perš, J., Muhovič, J., Košir, A., \& Murovec, B. (2026). \textit{Grading Handwritten Engineering Exams with Multimodal Large Language Models}. arXiv:2601.00730v1 [cs.CV], 2. januar 2026.}
  }
}

\NoTextSlide{
  title = {Arhitektura: ne en prompt, ampak cevovod},
  subtitle = {GitHub repo še ni javen, na željo lahko dobijo dostop zainteresirani pedagogi.},
  body = {
    \placeimage{72}{130}{820}{316}{stem-pipeline-overview-nocaption.png}
    \SlaifFootnotePanel[fill=SlaifPanelBlue]{Perš, J., Muhovič, J., Košir, A., \& Murovec, B. (2026). \textit{Grading Handwritten Engineering Exams with Multimodal Large Language Models}. arXiv:2601.00730v1 [cs.CV], 2. januar 2026.}
  }
}

\NoTextSlide{
  title = {Kaj je novost in zakaj je evalvacija težka},
  subtitle = {Pomemben rezultat ni "model zna ocenjevati", ampak sistemska zasnova okoli modela.},
  body = {
    \ColoredPanels{
      \ColoredPanel{Novost}{}{\SlaifPanelCenteredItems{\SlaifPanelTextItem{realistični skeni STEM rešitev\\brez prilagajanja študentov\\brez ročnega pretipkavanja}}}{}
      \ColoredPanel{Sistemska zasnova}{}{\SlaifPanelCenteredItems{\SlaifPanelTextItem{referenčna rešitev\\pravila in stroge predloge\\več neodvisnih ocen\\nadzorni korak}}}{}
      \ColoredPanel{Težka evalvacija}{}{\SlaifPanelCenteredItems{\SlaifPanelTextItem{zasebnost podatkov\\redki javni nabori\\človeška ocena ni absolutna\\pomembna je tudi razlaga}}}{}
    }
    \SlaifPlainBulletBlock[variant=blue,x=90,y=380,width=790]{
      \item Poleg točnosti je pomembna tudi nepristranskost ocenjevanja.
      \item Rezultati morajo biti stabilni med ponovitvami.
      \item Sistem mora zaznati primere, ki zahtevajo ročni pregled.
      \item Povratna informacija mora študentu razložiti, zakaj je rešitev pravilna, delno pravilna ali napačna.
    }
  }
}

\NoTextSlide{
  title = {Referenčna rešitev: zakaj je ključna},
  subtitle = {Ocenjevanje ni "model proti študentu", ampak model ob pravilih in lokalni normi.},
  body = {
    \ColoredPanels{
      \ColoredPanel{1. Referenca}{}{\SlaifPanelCenteredItems{\SlaifPanelTextItem{Predavatelj pripravi 100\% rešitev, lahko tudi ročno pisano.}}}{}
      \ColoredPanel{2. Povzetek}{}{\SlaifPanelCenteredItems{\SlaifPanelTextItem{Sistem pretvori referenco v ocenjevalno sidro.}}}{}
      \ColoredPanel{3. Ocenjevanje}{}{\SlaifPanelCenteredItems{\SlaifPanelTextItem{Odgovor se presoja ob pravilih, ne ob splošnem vtisu modela.}}}{}
    }
    \SlaifCommentaryPanel{Zakaj referenca ni dodatek}{Brez reference se model lažje premakne v splošno razlago in sistematično precenjevanje.}
  }
}

\NoTextSlide{
  title = {Prisotnost odgovora in stroge predloge},
  subtitle = {Dve majhni varovali, ki preprečita velik del nereda.},
  body = {
    \ColoredPanels{
      \ColoredPanel{Prisotnost odgovora}{}{\SlaifPanelCenteredItems{\SlaifPanelTextItem{Najprej preveri, ali ima vprašanje dejanski odgovor. Praznine se ne smejo spremeniti v delne točke.}}}{}
      \ColoredPanel{Stroga predloga}{}{\SlaifPanelCenteredItems{\SlaifPanelTextItem{Poročilo ima pričakovano strukturo: vprašanje, povzetek, pravila, razlaga in točke.}}}{}
      \ColoredPanel{Validacija izvoza}{}{\SlaifPanelCenteredItems{\SlaifPanelTextItem{Če vsota ne ustreza prispevkom ali format ni strojno berljiv, sledi ročni pregled.}}}{}
    }
    \SlaifCommentaryPanel{Zakaj stroga oblika šteje}{Cevovod zaradi svoje zasnove izdela strukturiran zapis, ki ga lahko deterministično razčlenimo, preverimo, preračunamo in prenesemo v tabele ali poročila.}
  }
}

\NoTextSlide{
  title = {Zakaj več ocenjevalcev in nadzornik?},
  subtitle = {En sam modelni odgovor je premalo stabilen za pedagoško rabo.},
  body = {
    \ColoredPanels{
      \ColoredPanel{Ocenjevalec A}{}{\SlaifPanelCenteredItems{\SlaifPanelTextItem{neodvisen osnutek}}}{}
      \ColoredPanel{Ocenjevalec B}{}{\SlaifPanelCenteredItems{\SlaifPanelTextItem{neodvisen osnutek}}}{}
      \ColoredPanel{Ocenjevalec C}{}{\SlaifPanelCenteredItems{\SlaifPanelTextItem{neodvisen osnutek}}}{}
      \ColoredPanel{Nadzornik}{}{\SlaifPanelCenteredItems{\SlaifPanelTextItem{združi osnutke, uveljavi predlogo in označi neskladja}}}{}
    }
    \SlaifCommentaryPanel{Interpretacija}{Nesoglasje ni napaka; je signal za človeški pregled.}
  }
}

\TwoColumnNoTextSlide{
  title = {Primer referenčne rešitve},
  subtitle = {Referenčna rešitev model zasidra na pravo mesto glede pričakovanega nivoja znanja in podrobnosti.},
  left = {
    \SlaifPlainBulletBlock[variant=blue,x=72,y=178,width=390,height=260]{
      \item ni nujno tipkana; lahko je ročno pisana;
      \item vsebuje metode, skice in formule;
      \item služi kot lokalna ocenjevalna norma;
      \item mora biti dovolj jasna, da jo sistem povzame.
    }
  },
  right = {
    \placeimage{630}{118}{223}{300}{supplement-reference-solution-page1.png}
    \SlaifFootnotePanel[fill=SlaifPanelBlue]{Perš, J., Muhovič, J., Košir, A., \& Murovec, B. (2026). \textit{Grading Handwritten Engineering Exams with Multimodal Large Language Models}. arXiv:2601.00730v1 [cs.CV], 2. januar 2026.}
  }
}

\NoTextSlide{
  title = {Izhod ni samo številka, ampak obsežen dokument},
  subtitle = {Glavna prednost za študente: obsežna povratna informacija s podrobno razlago, kaj so naredili narobe.},
  body = {
    \placeimage{54}{126}{410}{312}{supplement-final-grade.png}
    \placeimage{510}{126}{377}{312}{supplement-task-feedback.png}
    \SlaifFootnotePanel[fill=SlaifPanelBlue]{Perš, J., Muhovič, J., Košir, A., \& Murovec, B. (2026). \textit{Grading Handwritten Engineering Exams with Multimodal Large Language Models}. arXiv:2601.00730v1 [cs.CV], 2. januar 2026.}
  }
}

\NoTextSlide{
  title = {Povratna informacija od študentov},
  subtitle = {Predmet: Elektronika z digitalno tehniko, zimski semester 2025/2026},
  body = {
    \ColoredPanels{
      \ColoredPanel{Vprašanje 1}{Odnos do UI-prvega ocenjevanja}{\SlaifSliderTable{
        labelX=54,y=265,labelWidth=150,barX=222,barWidth=70,barHeight=12,rowStep=28,
        valueYOffset=2.8,valueWidth=42,maxValue=100,valueSuffix={\%},
        rows={pozitiven odnos ali preferenca/64,negativen odnos/21,neopredeljeni/14}
      }}{}
      \ColoredPanel{Vprašanje 2}{Ali imate korist od sistema?}{\SlaifSliderTable{
        labelX=370,y=265,labelWidth=62,barX=450,barWidth=108,barHeight=12,rowStep=32,
        valueYOffset=2.8,valueWidth=42,maxValue=100,valueSuffix={\%},
        rows={da/71,ne/29}
      }}{}
      \ColoredPanel{Vprašanje 3}{Katere koristi vidite?}{\SlaifSliderTable{
        labelX=664,y=265,labelWidth=122,barX=804,barWidth=62,barHeight=12,rowStep=28,
        valueYOffset=2.8,valueWidth=42,maxValue=100,valueSuffix={\%},
        rows={podrobne razlage/41,občutek nepristranskosti/34,hiter odziv/21}
      }}{}
    }
    \SlaifFootnotePanel[fill=SlaifPanelBlue]{Perš, J., Muhovič, J., Košir, A., \& Murovec, B. (2026). \textit{Grading Handwritten Engineering Exams with Multimodal Large Language Models}. arXiv:2601.00730v1 [cs.CV], 2. januar 2026. V anketi je sodelovalo 14 študentov po 8 tedenskih UI-ocenjenih kvizih.}
  }
}

\NoTextSlide{
  title = {Povratna informacija od študentov},
  subtitle = {Predmet: Elektronika z digitalno tehniko, zimski semester 2025/2026},
  body = {
    \ColoredPanels{
      \ColoredPanel{Vprašanje 4}{Kaj vas skrbi?}{\SlaifSliderTable{
        labelX=62,y=252,labelWidth=205,barX=292,barWidth=145,barHeight=13,rowStep=34,
        valueYOffset=3.0,valueWidth=48,maxValue=100,valueSuffix={\%},
        rows={spregledani odgovori/35,več napak kot pri profesorjih/26,izpit je manj zahteven/22,izpit je bolj zahteven/17}
      }}{}
      \ColoredPanel{Vprašanje 5}{Skupna presoja}{\SlaifSliderTable{
        labelX=530,y=252,labelWidth=215,barX=770,barWidth=100,barHeight=13,rowStep=36,
        valueYOffset=3.0,valueWidth=48,maxValue=100,valueSuffix={\%},
        rows={prednosti pretehtajo slabosti/43,slabosti pretehtajo prednosti/29,brez bistvene razlike/21}
      }}{}
    }
    \SlaifFootnotePanel[fill=SlaifPanelBlue]{Perš, J., Muhovič, J., Košir, A., \& Murovec, B. (2026). \textit{Grading Handwritten Engineering Exams with Multimodal Large Language Models}. arXiv:2601.00730v1 [cs.CV], 2. januar 2026. V anketi je sodelovalo 14 študentov po 8 tedenskih UI-ocenjenih kvizih.}
  }
}

\TwoColumnNoTextSlide{
  title = {Kaj povedo rezultati},
  subtitle = {Najboljši sodobni modeli so uporabni, šibkejši pa hitro odpovejo.},
  left = {
    \SlaifPlainBulletBlock[variant=blue,x=72,y=178,width=390,height=260]{
      \item MAD: povprečna absolutna razlika do učiteljeve ocene;
      \item STD(|\(\Delta\)|): razpršenost absolutnih napak med študenti;
      \item bias: sistematično precenjevanje ali podcenjevanje;
      \item najboljši sistemi dosežejo enomestni MAD z nizkim biasom.
    }
  },
  right = {
    \placeimage{575}{95}{325}{345}{stem-baseline-metrics-crop.png}
    \SlaifFootnotePanel[fill=SlaifPanelBlue]{Perš, J., Muhovič, J., Košir, A., \& Murovec, B. (2026). \textit{Grading Handwritten Engineering Exams with Multimodal Large Language Models}. arXiv:2601.00730v1 [cs.CV], 2. januar 2026.}
  }
}

\TwoColumnNoTextSlide{
  title = {Zakaj "samo vprašaj model" ni dovolj},
  subtitle = {Ablacija: trivialno pozivanje brez pravil, reference in nadzornika poslabša ocenjevanje.},
  left = {
    \SlaifPlainBulletBlock[variant=blue,x=72,y=178,width=390,height=260]{
      \item referenca ni dodatek, ampak sidro;
      \item pravila morajo biti vidna v razlagi;
      \item predloga prepreči formatni kaos;
      \item nadzornik in validacija zmanjšata zdrse.
    }
  },
  right = {
    \placeimage{535}{95}{371}{345}{stem-guidance-ablation-crop.png}
    \SlaifFootnotePanel[fill=SlaifPanelBlue]{Perš, J., Muhovič, J., Košir, A., \& Murovec, B. (2026). \textit{Grading Handwritten Engineering Exams with Multimodal Large Language Models}. arXiv:2601.00730v1 [cs.CV], 2. januar 2026.}
  }
}

\TwoColumnNoTextSlide{
  title = {Ročni pregled kot varnostni ventil},
  subtitle = {Ne avtomatiziramo vsega; avtomatiziramo pripravo in označimo sporne primere.},
  left = {
    \SlaifPlainBulletBlock[variant=blue,x=72,y=178,width=390,height=260]{
      \item kadar se ocenjevalci ne strinjajo, naj človek pogleda;
      \item prag ni tehnična malenkost, ampak pedagoška odločitev;
      \item cilj je zmanjšati obseg pregleda, ne odstraniti odgovornost;
      \item tako dobimo nadzorovan kompromis med hitrostjo in varnostjo.
    }
  },
  right = {
    \placeimage{510}{86}{396}{366}{stem-manual-review-triggers-crop.png}
    \SlaifFootnotePanel[fill=SlaifPanelBlue]{Perš, J., Muhovič, J., Košir, A., \& Murovec, B. (2026). \textit{Grading Handwritten Engineering Exams with Multimodal Large Language Models}. arXiv:2601.00730v1 [cs.CV], 2. januar 2026.}
  }
}

\NoTextSlide{
  title = {Etika in lokalna pravila},
  subtitle = {Sistem mora biti tak, da ga predavatelj razume dovolj dobro, da ga lahko zavrne.},
  body = {
    \ColoredPanels{
      \ColoredPanel{Podatki}{}{\SlaifPanelCenteredItems{\SlaifPanelTextItem{psevdoanonimizacija, lokalna deanonimizacija in jasen nadzor nad vhodnimi datotekami}}}{}
      \ColoredPanel{Študent}{}{\SlaifPanelCenteredItems{\SlaifPanelTextItem{ve, kateri del postopka je avtomatiziran in kje odloča človek}}}{}
      \ColoredPanel{Napake}{}{\SlaifPanelCenteredItems{\SlaifPanelTextItem{sistem pričakuje napačno branje, spregledan diagram in preveč samozavestno razlago}}}{}
    }
    \SlaifCommentaryPanel{Minimalni standard}{V izobraževanju ni dovolj, da je ocena približno pravilna; postopek mora biti razložljiv in popravljiv.}
  }
}

\NoTextSlide{
  title = {Kaj bi bil dober naslednji korak pri nas?},
  subtitle = {Ne velika obljuba, ampak majhen, preverljiv pilot.},
  body = {
    \ColoredPanels{
      \ColoredPanel{1. Izberemo}{}{\SlaifPanelCenteredItems{\SlaifPanelTextItem{en predmet, eno vrsto naloge in delovni tok, ki res obremenjuje ekipo}}}{}
      \ColoredPanel{2. Zberemo}{}{\SlaifPanelCenteredItems{\SlaifPanelTextItem{navodila, primere oddaj, kriterije in stare povratne informacije}}}{}
      \ColoredPanel{3. Naredimo}{}{\SlaifPanelCenteredItems{\SlaifPanelTextItem{osnutek prosojnic, vprašanj za zagovor ali ocenjevalnega poročila}}}{}
      \ColoredPanel{4. Merimo}{}{\SlaifPanelCenteredItems{\SlaifPanelTextItem{čas\\napake\\popravljivost\\nadzor}}}{}
    }
    \SlaifCommentaryPanel{Naslednji korak}{Cilj jesenske delavnice: konkretna opravila iz vaših predmetov, ne abstraktna razprava o UI.}
  }
}

\FinalSlide[slaif-si-qr-rgb.png]{Hvala za pozornost!}

\end{document}
